\documentclass[../LuanVan.tex]{subfiles}
\begin{document}

\section{Tổng quan hệ thống}
\label{section:4.1}

Hệ thống được xây dựng cho một cửa hàng, phục vụ ba nhóm công việc chính: quản lý dữ liệu CRM, ghi nhận hành trình tại các điểm chạm và trình bày kết quả nhận dạng biểu cảm.

Các chức năng đã xây dựng gồm:
\begin{itemize}
    \item quản lý khách hàng, sản phẩm, đơn hàng và lịch sử mua hàng;
    \item quản lý điểm chạm, lượt ghé thăm và các quan sát trong từng lượt;
    \item đăng ký mẫu khuôn mặt và tìm khách hàng bằng ảnh;
    \item hiển thị hành trình của từng khách hàng;
    \item thống kê nhãn biểu cảm theo điểm chạm, theo thời gian và theo khu vực;
    \item thống kê sự thay đổi nhãn giữa các điểm chạm và chất lượng dữ liệu đầu vào.
\end{itemize}

\subsection{Kiến trúc hệ thống}
\label{subsection:4.1.1}

Hệ thống gồm bốn thành phần chính chạy độc lập. Vai trò và công nghệ của từng thành phần được trình bày tại Bảng~\ref{tab:system_components}. Dịch vụ cập nhật cấu trúc cơ sở dữ liệu chỉ chạy khi khởi động và kết thúc sau khi hoàn thành.

\begin{table}[H]
    \centering
    \caption{Các thành phần đã triển khai}
    \label{tab:system_components}
    \resizebox{\textwidth}{!}{%
    \begin{tabular}{|p{2.5cm}|p{4.0cm}|p{8.0cm}|}
        \hline
        \textbf{Thành phần} & \textbf{Công nghệ} & \textbf{Chức năng đã triển khai} \\
        \hline
        \texttt{web} & React, TypeScript, Ant Design, ECharts & Giao diện quản lý khách hàng, sản phẩm, đơn hàng, điểm chạm, hành trình và báo cáo. \\
        \hline
        \texttt{api} & FastAPI, SQLAlchemy, Alembic & Xử lý nghiệp vụ, xác thực, liên kết lượt ghé thăm, tổng hợp báo cáo và cung cấp API. \\
        \hline
        \texttt{vision} & OpenCV; mã tích hợp RetinaFace, DeepFace và ArcFace & Kiểm tra ảnh, phát hiện khuôn mặt, phân loại biểu cảm và tạo véc-tơ khuôn mặt theo cấu hình chạy. \\
        \hline
        \texttt{postgres} & PostgreSQL 16, pgvector 0.8.3 & Lưu hồ sơ CRM, lượt ghé thăm, quan sát, đơn hàng và véc-tơ khuôn mặt. \\
        \hline
    \end{tabular}}
\end{table}

Hình~\ref{fig:system_architecture_ch4} thể hiện quan hệ giữa các thành phần. Giao diện CRM không trực tiếp thực hiện tính năng xử lý ảnh hoặc truy cập cơ sở dữ liệu. Dữ liệu từ nguồn ảnh đi qua máy chủ nghiệp vụ trước khi được xử lý và lưu trữ.

\begin{figure}[H]
    \centering
    \resizebox{0.98\textwidth}{!}{%
    \begin{tikzpicture}[
        node distance=9mm,
        block/.style={draw, rounded corners=2pt, align=center, minimum height=12mm, text width=3.0cm, font=\small},
        store/.style={draw, cylinder, shape border rotate=90, aspect=0.25, align=center, minimum height=14mm, text width=2.8cm, font=\small},
        flow/.style={-{Latex[length=2mm]}, thick}
    ]
        \node[block] (browser) {Trình duyệt\\Giao diện CRM};
        \node[block, below=of browser] (camera) {Nguồn ảnh\\tại điểm chạm};
        \node[block, right=18mm of camera] (api) {Máy chủ\\nghiệp vụ};
        \node[block, right=18mm of api, yshift=10mm] (vision) {Tính năng\\xử lý ảnh};
        \node[store, right=18mm of api, yshift=-13mm] (db) {PostgreSQL\\và pgvector};
        \draw[flow, <->] (browser) -- (api);
        \draw[flow] (camera) -- (api);
        \draw[flow, <->] (api) -- (vision);
        \draw[flow, <->] (api) -- (db);
    \end{tikzpicture}%
    }
    \caption{Kiến trúc của hệ thống đã xây dựng}
    \label{fig:system_architecture_ch4}
\end{figure}

\subsection{Luồng dữ liệu chính}
\label{subsection:4.1.2}

Một yêu cầu gồm ảnh, mã điểm chạm và thời gian ghi nhận được gửi đến máy chủ nghiệp vụ. Máy chủ kiểm tra dữ liệu rồi chuyển ảnh tới tính năng xử lý ảnh. Kết quả biểu cảm và véc-tơ khuôn mặt được trả về để máy chủ đối sánh khách hàng, tạo hoặc cập nhật lượt ghé thăm và lưu bản ghi quan sát. Sau khi lưu, dữ liệu được dùng cho màn hình hành trình, bảng phân bố, biểu đồ thay đổi nhãn và báo cáo chất lượng dữ liệu.

\begin{figure}[H]
    \centering
    \resizebox{0.96\textwidth}{!}{%
    \begin{tikzpicture}[
        node distance=7mm,
        block/.style={draw, rounded corners=2pt, align=center, minimum height=11mm, text width=2.8cm, font=\small},
        flow/.style={-{Latex[length=2mm]}, thick}
    ]
        \node[block] (source) {Ảnh tại\\điểm chạm};
        \node[block, right=of source] (process) {Kiểm tra và\\xử lý dữ liệu};
        \node[block, right=of process] (visit) {Tạo hoặc cập nhật\\lượt ghé thăm};
        \node[block, right=of visit] (store) {Lưu khách hàng,\\quan sát và đơn hàng};
        \node[block, right=of store] (display) {Tra cứu và\\trình bày báo cáo};
        \draw[flow] (source) -- (process);
        \draw[flow] (process) -- (visit);
        \draw[flow] (visit) -- (store);
        \draw[flow] (store) -- (display);
    \end{tikzpicture}%
    }
    \caption{Luồng dữ liệu chính trong hệ thống}
    \label{fig:system_data_flow_ch4}
\end{figure}

\section{Cài đặt hệ thống}
\label{section:4.2}

\subsection{Ca sử dụng tổng quan}
\label{subsection:4.2.1}

Hệ thống có một tác nhân chính là người quản lý cửa hàng. Sau khi đăng nhập, người quản lý có thể thao tác với dữ liệu CRM, theo dõi hành trình và xem các báo cáo biểu cảm.

\begin{figure}[H]
    \centering
    \begin{tikzpicture}[
        usecase/.style={draw, ellipse, align=center, minimum width=4.6cm, minimum height=10mm, font=\small},
        actor/.style={draw, rounded corners=2pt, align=center, minimum width=3.2cm, minimum height=12mm, font=\small},
        link/.style={thick}
    ]
        \node[actor] (manager) {Người quản lý\\cửa hàng};
        \node[usecase, right=24mm of manager, yshift=25mm] (crm) {Quản lý khách hàng, sản phẩm\\và lịch sử mua hàng};
        \node[usecase, right=24mm of manager, yshift=8mm] (journey) {Theo dõi hành trình\\tại các điểm chạm};
        \node[usecase, right=24mm of manager, yshift=-9mm] (reports) {Xem báo cáo phân bố và\\thay đổi nhãn biểu cảm};
        \node[usecase, right=24mm of manager, yshift=-26mm] (timeline) {Theo dõi biến thiên theo\\thời gian và khu vực};
        \draw[link] (manager) -- (crm);
        \draw[link] (manager) -- (journey);
        \draw[link] (manager) -- (reports);
        \draw[link] (manager) -- (timeline);
    \end{tikzpicture}
    \caption{Ca sử dụng tổng quan của hệ thống}
    \label{fig:overall_usecase_ch4}
\end{figure}

\clearpage
\subsection{Phân rã chức năng}
\label{subsection:4.2.2}

Các chức năng được chia thành ba nhóm như Hình~\ref{fig:function_decomposition_ch4}. Nhóm CRM quản lý dữ liệu nghiệp vụ. Nhóm hành trình quản lý điểm chạm và lượt ghé thăm. Nhóm phân tích đọc các bản ghi hợp lệ để tạo báo cáo.

\begin{figure}[H]
    \centering
    \resizebox{\textwidth}{!}{%
    \begin{tikzpicture}[
        node distance=7mm,
        root/.style={draw, rounded corners=2pt, align=center, minimum height=11mm, text width=3.0cm, font=\small\bfseries},
        group/.style={draw, rounded corners=2pt, align=center, minimum height=11mm, text width=3.0cm, font=\small},
        item/.style={draw, align=center, minimum height=10mm, text width=3.0cm, font=\scriptsize},
        flow/.style={-{Latex[length=2mm]}, thick}
    ]
        \node[root] (system) {Touchpoint CRM};
        \node[group, right=12mm of system, yshift=20mm] (crm) {Quản lý CRM};
        \node[group, right=12mm of system] (journey) {Quản lý hành trình};
        \node[group, right=12mm of system, yshift=-20mm] (analysis) {Phân tích biểu cảm};
        \node[item, right=7mm of crm] (crmitems) {Khách hàng\\Sản phẩm\\Đơn hàng};
        \node[item, right=7mm of journey] (journeyitems) {Điểm chạm\\Lượt ghé thăm\\Quan sát};
        \node[item, right=7mm of analysis] (analysisitems) {Phân bố\\Biến thiên\\Dữ liệu lỗi};
        \draw[flow] (system) -- (crm);
        \draw[flow] (system) -- (journey);
        \draw[flow] (system) -- (analysis);
        \draw[flow] (crm) -- (crmitems);
        \draw[flow] (journey) -- (journeyitems);
        \draw[flow] (analysis) -- (analysisitems);
    \end{tikzpicture}%
    }
    \caption{Phân rã các nhóm chức năng của hệ thống}
    \label{fig:function_decomposition_ch4}
\end{figure}

\subsection{Đặc tả các chức năng chính}
\label{subsection:4.2.3}

Bảng~\ref{tab:main_usecases_ch4} mô tả các chức năng chính của hệ thống. Các thao tác thêm, sửa và tra cứu thông thường được gộp theo nhóm để tránh lặp lại nội dung.

\begin{table}[H]
    \centering
    \caption{Đặc tả các chức năng chính}
    \label{tab:main_usecases_ch4}
    \resizebox{\textwidth}{!}{%
    \begin{tabular}{|p{1.5cm}|p{3.7cm}|p{4.8cm}|p{5.0cm}|}
        \hline
        \textbf{Mã} & \textbf{Chức năng} & \textbf{Điều kiện và dữ liệu vào} & \textbf{Kết quả} \\
        \hline
        UC01 & Tra cứu khách hàng và lịch sử mua hàng & Người quản lý đã đăng nhập; mã, tên hoặc số điện thoại khách hàng & Hồ sơ, ảnh đại diện và các đơn hàng của đúng khách hàng \\
        \hline
        UC02 & Ghi nhận quan sát tại điểm chạm & Có ảnh, mã điểm chạm và thời gian ghi nhận & Trạng thái xử lý ảnh, khách hàng được xác định và bản ghi quan sát \\
        \hline
        UC03 & Theo dõi một lượt ghé thăm & Chọn một lượt đã có bản ghi quan sát & Chuỗi điểm chạm được sắp xếp theo thời gian cùng nhãn biểu cảm \\
        \hline
        UC04 & Xem biến thiên theo thời gian và khu vực & Chọn điểm chạm, ngày và độ dài khoảng thời gian & Biểu đồ số quan sát của bảy nhãn trong cùng khu vực theo từng mốc thời gian \\
        \hline
        UC05 & Xem báo cáo chất lượng dữ liệu & Có dữ liệu quan sát hợp lệ và lỗi & Số bản ghi hợp lệ, không có khuôn mặt hoặc ảnh không hợp lệ \\
        \hline
    \end{tabular}}
\end{table}

Đối với UC04, máy chủ chỉ lấy các quan sát hợp lệ tại điểm chạm đã chọn, chia thời gian thành các khoảng 5, 15, 30 hoặc 60 phút rồi đếm từng nhãn trong mỗi khoảng. Người dùng chọn thêm ngày để các mốc giờ trên biểu đồ thuộc cùng một ngày và cùng một khu vực. Chức năng không quy đổi nhãn thành điểm tích cực--tiêu cực.

\subsection{Xây dựng API và tính năng xử lý ảnh}
\label{subsection:4.2.4}

Máy chủ cung cấp các đường dẫn quản lý khách hàng, sản phẩm, đơn hàng, điểm chạm, lượt ghé thăm, quan sát và báo cáo. Tài liệu OpenAPI tại Hình~\ref{fig:api_document_ch4} được tạo trực tiếp từ mã nguồn đang chạy. Mã sự kiện của một quan sát có ràng buộc duy nhất để một ảnh gửi lại không tạo hai bản ghi.

\begin{figure}[H]
    \centering
    \includegraphics[width=0.98\textwidth]{Hinhve/Chuong4/4_11_tai_lieu_api.png}
    \caption{Các API chính của máy chủ nghiệp vụ}
    \label{fig:api_document_ch4}
\end{figure}

Tính năng xử lý ảnh trả về trạng thái ảnh, kết quả biểu cảm và véc-tơ đặc trưng khuôn mặt. Máy chủ dùng véc-tơ để đối sánh với mẫu đã đăng ký rồi mới quyết định mã khách hàng. Tính năng này không truy cập hồ sơ CRM và không tự gán danh tính. Hình~\ref{fig:image_pipeline_ch4} trình bày luồng đã cài đặt.

\begin{figure}[H]
    \centering
    \resizebox{0.96\textwidth}{!}{%
    \begin{tikzpicture}[
        node distance=7mm,
        block/.style={draw, rounded corners=2pt, align=center, minimum height=11mm, text width=2.8cm, font=\small},
        decision/.style={draw, diamond, aspect=2.0, align=center, text width=2.3cm, font=\small},
        flow/.style={-{Latex[length=2mm]}, thick}
    ]
        \node[block] (input) {Ảnh và thông tin\\điểm chạm};
        \node[block, right=of input] (decode) {Giải mã và\\kiểm tra ảnh};
        \node[block, right=of decode] (detect) {Phát hiện\\khuôn mặt};
        \node[decision, right=of detect] (valid) {Một khuôn mặt\\hợp lệ?};
        \node[block, right=of valid, yshift=9mm] (tasks) {Phân loại biểu cảm\\và tạo đặc trưng};
        \node[block, right=of valid, yshift=-9mm] (error) {Lưu trạng thái lỗi};
        \node[block, right=of tasks, yshift=-9mm] (record) {Tạo bản ghi\\quan sát};
        \draw[flow] (input) -- (decode);
        \draw[flow] (decode) -- (detect);
        \draw[flow] (detect) -- (valid);
        \draw[flow] (valid) -- node[above, font=\scriptsize] {Có} (tasks);
        \draw[flow] (valid) -- node[below, font=\scriptsize] {Không} (error);
        \draw[flow] (tasks) -- (record);
        \draw[flow] (error) -- (record);
    \end{tikzpicture}%
    }
    \caption{Luồng tiếp nhận và xử lý một ảnh tại điểm chạm}
    \label{fig:image_pipeline_ch4}
\end{figure}

Mã nguồn có cấu hình nhẹ để kiểm tra giao tiếp và trạng thái lỗi, cùng cấu hình tích hợp RetinaFace, mô hình Emotion của DeepFace và ArcFace. Tại thời điểm ghi nhận Chương 4, hệ thống Docker đang chạy cấu hình nhẹ. Cấu hình mô hình thật, trọng số và ngưỡng đối sánh chưa được kiểm chứng đầy đủ; nội dung đánh giá được trình bày riêng tại Chương 5.

\section{Giao diện CRM}
\label{section:4.3}

\subsection{Màn hình tổng quan}
\label{subsection:4.3.1}

Trang tổng quan lấy số liệu trực tiếp từ API và trình bày số khách hàng, sản phẩm, lượt ghé thăm cùng phân bố nhãn theo điểm chạm. Giao diện được minh họa tại Hình~\ref{fig:dashboard_ch4}.

\begin{figure}[H]
    \centering
    \includegraphics[width=0.98\textwidth]{Hinhve/Chuong4/4_01_tong_quan.png}
    \caption{Trang tổng quan CRM}
    \label{fig:dashboard_ch4}
\end{figure}

Tab thứ 2 trên màn hình tổng quan dùng để theo dõi biến thiên trong cùng một khu vực và cùng một ngày. Người dùng chọn điểm chạm, ngày và khoảng gom thời gian. Hình~\ref{fig:timeline_area_ch4} minh họa cách hệ thống hiển thị số quan sát của bảy nhãn theo từng khoảng thời gian.

\begin{figure}[H]
    \centering
    \includegraphics[width=0.98\textwidth]{Hinhve/Chuong4/4_13_bien_thien_theo_thoi_gian_khu_vuc.png}
    \caption{Theo dõi biến thiên nhãn theo thời gian trong cùng khu vực}
    \label{fig:timeline_area_ch4}
\end{figure}

\subsection{Màn hình khách hàng và lịch sử mua hàng}
\label{subsection:4.3.2}

Màn hình khách hàng hỗ trợ tìm theo mã, tên hoặc số điện thoại. Hồ sơ chi tiết hiển thị ảnh đại diện, thông tin liên hệ và các đơn hàng của riêng khách hàng. Hình~\ref{fig:customer_history_ch4} minh họa việc tra cứu hồ sơ và lịch sử mua hàng trên cùng một màn hình.

\begin{figure}[H]
    \centering
    \includegraphics[width=0.98\textwidth]{Hinhve/Chuong4/4_03_ho_so_lich_su_mua_hang.png}
    \caption{Hồ sơ khách hàng và lịch sử mua hàng}
    \label{fig:customer_history_ch4}
\end{figure}

\subsection{Màn hình sản phẩm và đơn hàng}
\label{subsection:4.3.3}

Mỗi sản phẩm có mã, tên, nhóm, giá hiện tại và trạng thái. Khi tạo đơn hàng, máy chủ sao chép mã, tên và đơn giá vào dòng đơn. Vì vậy, thay đổi giá hiện tại không làm thay đổi lịch sử đã phát sinh. Quy tắc này được kiểm thử tự động tại tầng API.

Hình~\ref{fig:products_ch4} cho thấy danh mục sản phẩm. Hình~\ref{fig:orders_ch4} cho thấy danh sách đơn hàng gồm mã đơn, thời gian, tổng tiền và trạng thái.

\begin{figure}[H]
    \centering
    \includegraphics[width=0.98\textwidth]{Hinhve/Chuong4/4_04_danh_muc_san_pham.png}
    \caption{Danh mục sản phẩm trong hệ thống CRM}
    \label{fig:products_ch4}
\end{figure}

\begin{figure}[H]
    \centering
    \includegraphics[width=0.98\textwidth]{Hinhve/Chuong4/4_05_lich_su_mua_hang.png}
    \caption{Danh sách đơn hàng trong hệ thống CRM}
    \label{fig:orders_ch4}
\end{figure}

\subsection{Màn hình điểm chạm và hành trình}
\label{subsection:4.3.4}

Hệ thống dùng bốn điểm chạm trong một cửa hàng: cửa vào, khu trưng bày sản phẩm, khu tư vấn và quầy thanh toán. Thứ tự được lưu trong cơ sở dữ liệu và dùng khi hiển thị hành trình. Hình~\ref{fig:touchpoints_ch4} xác nhận cả bốn điểm chạm đang hoạt động.

\begin{figure}[H]
    \centering
    \includegraphics[width=0.98\textwidth]{Hinhve/Chuong4/4_06_diem_cham.png}
    \caption{Danh sách điểm chạm của cửa hàng}
    \label{fig:touchpoints_ch4}
\end{figure}

Màn hình theo dõi tại Hình~\ref{fig:journey_screen_ch4} hiển thị danh sách lượt ghé thăm, số khách hàng có hành trình, tổng số lượt và sơ đồ của lượt đang chọn. Mỗi nút gồm tên điểm chạm và nhãn biểu cảm; mũi tên thể hiện thứ tự thời gian. Màu chỉ dùng để phân biệt nhãn, không biểu thị mức tích cực hoặc tiêu cực.

\begin{figure}[H]
    \centering
    \includegraphics[width=0.98\textwidth]{Hinhve/Chuong4/4_07_theo_doi_hanh_trinh.png}
    \caption{Màn hình theo dõi hành trình của từng khách hàng}
    \label{fig:journey_screen_ch4}
\end{figure}

\subsection{Màn hình phân tích biểu cảm và chất lượng dữ liệu}
\label{subsection:4.3.5}

Máy chủ đếm các quan sát có trạng thái ảnh và biểu cảm hợp lệ theo từng điểm chạm, từng nhãn. Giao diện trình bày cả số lượng và tỷ lệ như Hình~\ref{fig:distribution_screen_ch4}. Các cột chỉ mô tả nhãn đã lưu, không được dùng để kết luận khách hàng hài lòng hay xác định nguyên nhân của biểu cảm.

\begin{figure}[H]
    \centering
    \includegraphics[width=0.98\textwidth]{Hinhve/Chuong4/4_08_phan_bo_bieu_cam.png}
    \caption{Phân bố nhãn biểu cảm tại bốn điểm chạm}
    \label{fig:distribution_screen_ch4}
\end{figure}

Để mô tả thay đổi giữa hai điểm chạm liên tiếp, máy chủ chọn các bản ghi trong cùng lượt ghé thăm, sắp xếp theo thời gian và đếm từng cặp nhãn trước--sau. Hình~\ref{fig:change_screen_ch4} thể hiện các cặp bằng biểu đồ luồng và giữ số đếm cụ thể trong bảng phía dưới. Đây không phải ma trận xác suất và không áp dụng mô hình Markov.

\begin{figure}[H]
    \centering
    \includegraphics[width=0.98\textwidth]{Hinhve/Chuong4/4_09_thay_doi_bieu_cam.png}
    \caption{Luồng thay đổi nhãn giữa các điểm chạm liên tiếp}
    \label{fig:change_screen_ch4}
\end{figure}

Màn hình chất lượng dữ liệu phân tách các quan sát hợp lệ, ảnh không có khuôn mặt và ảnh không hợp lệ. Cách trình bày các trạng thái xử lý được minh họa tại Hình~\ref{fig:data_quality_ch4}.

\begin{figure}[H]
    \centering
    \includegraphics[width=0.98\textwidth]{Hinhve/Chuong4/4_10_chat_luong_du_lieu.png}
    \caption{Báo cáo trạng thái xử lý dữ liệu đầu vào}
    \label{fig:data_quality_ch4}
\end{figure}

\section{Kiểm tra phiên bản triển khai}
\label{section:4.4}

\subsection{Kết quả kiểm thử}
\label{subsection:4.4.1}

Bảng~\ref{tab:test_results_ch4} tổng hợp các kiểm tra kỹ thuật đã chạy. Sáu kiểm thử API gồm quy tắc giữ giá lịch sử, quản lý điểm chạm, ghi nhận quan sát và báo cáo theo thời gian tại một điểm chạm. Hai kiểm thử tính năng xử lý ảnh kiểm tra trạng thái xử lý của cấu hình nhẹ. Giao diện được biên dịch ở chế độ phát hành trước khi tạo ảnh chụp.

\begin{table}[H]
    \centering
    \caption{Kết quả kiểm tra phiên bản triển khai}
    \label{tab:test_results_ch4}
    \resizebox{0.9\textwidth}{!}{%
    \begin{tabular}{|p{4.8cm}|p{4.2cm}|p{4.0cm}|}
        \hline
        \textbf{Phần kiểm tra} & \textbf{Kết quả} & \textbf{Trạng thái} \\
        \hline
        Máy chủ nghiệp vụ và quy tắc CRM & 6/6 kiểm thử đạt & Đạt \\
        \hline
        Tính năng xử lý ảnh ở cấu hình nhẹ & 2/2 kiểm thử đạt & Đạt \\
        \hline
        Biên dịch React và TypeScript & Tạo gói phát hành thành công & Đạt \\
        \hline
        Di trú dữ liệu PostgreSQL & Cập nhật đến phiên bản mới nhất & Đạt \\
        \hline
    \end{tabular}}
\end{table}

\subsection{Giới hạn của phiên bản hiện tại}
\label{subsection:4.4.2}

Phiên bản hiện tại còn các giới hạn sau:
\begin{itemize}
    \item Chưa kết nối camera vật lý trong phiên bản đang kiểm tra.
    \item Cấu hình mô hình thật chưa được kiểm chứng trong môi trường Docker hiện tại; ngưỡng nhận dạng khách hàng và độ chính xác FER chưa được đánh giá.
    \item Hệ thống phục vụ một cửa hàng, chưa có chức năng quản lý nhiều chi nhánh.
    \item Báo cáo chỉ mô tả phân bố và sự thay đổi của nhãn biểu cảm quan sát được; hệ thống không suy luận trạng thái tâm lý, mức hài lòng hoặc nguyên nhân hành vi mua hàng.
\end{itemize}

Các giới hạn được tách khỏi kết quả đã kiểm tra. Chương này chỉ trình bày phần mềm và các luồng CRM đã xây dựng. Việc tạo dữ liệu khi chưa có camera và kết quả chạy thử được trình bày tại Chương 5.

\end{document}
